% source: RH-Framework/notebooks/T03_ACTION_DOMAIN_JUPYTER.md
\hypertarget{t03-action-domain-verification-jupyter-procedure}{%
\chapter{T03 Action-Domain Verification --- Jupyter Procedure}\label{t03-action-domain-verification-jupyter-procedure}}

This procedure uses Python cells only. No PowerShell is required.

\hypertarget{set-repository-paths}{%
\section{1. Set repository paths}\label{set-repository-paths}}

\begin{verbatim}
from pathlib import Path
import subprocess
import sys

RH_REPO = Path(
    r"C:\Users\abc\Desktop\New Folder (2)\New folder\PROVISNAL RELATED\RH-Framework"
).resolve()

MP_REPO = Path(
    r"C:\Users\abc\Desktop\New Folder (2)\New folder\PROVISNAL RELATED\MP_clock_free_endpoint_17e51005"
).resolve()

OPERATOR_COMMIT = "abc9725a8f232feab81ba63230085457913aadd7"
STATE_COMMIT = "f9e2f71c2d0ef7db239d7e12b18394eff8e83a1e"

OPERATOR_WORKTREE = MP_REPO.parent / "MP_operator_source_abc9725"
STATE_WORKTREE = MP_REPO.parent / "MP_state_source_f9e2f71"

T03_C1_JSON = (
    RH_REPO / "certificates" / "T03_C1_PINNED_ACTION_INVENTORY_v0.1.json"
)
T03_C1_NPZ = (
    RH_REPO / "certificates" / "T03_C1_PINNED_ACTION_ARRAYS_v0.1.npz"
)

print("RH framework:", RH_REPO)
print("MP source clone:", MP_REPO)
print("operator worktree:", OPERATOR_WORKTREE)
print("state worktree:", STATE_WORKTREE)
\end{verbatim}

Change only the local paths when your folders differ.

\hypertarget{git-helper-and-rh-branch-audit}{%
\section{2. Git helper and RH branch audit}\label{git-helper-and-rh-branch-audit}}

\begin{verbatim}
def run_git(repo: Path, *args: str) -> str:
    completed = subprocess.run(
        ["git", "-C", str(repo), *args],
        check=True,
        capture_output=True,
        text=True,
    )
    return completed.stdout.strip()

print("RH branch:", run_git(RH_REPO, "branch", "--show-current"))
print("RH head:", run_git(RH_REPO, "rev-parse", "HEAD"))
\end{verbatim}

Expected branch:

\begin{verbatim}
agent/t03-action-domain-adapter
\end{verbatim}

\hypertarget{install-the-rh-verification-package}{%
\section{3. Install the RH verification package}\label{install-the-rh-verification-package}}

\begin{verbatim}
completed = subprocess.run(
    [sys.executable, "-m", "pip", "install", "-e", str(RH_REPO)],
    text=True,
    capture_output=True,
)
print(completed.stdout)
print(completed.stderr)
print("return code:", completed.returncode)
assert completed.returncode == 0
\end{verbatim}

\hypertarget{verify-t03-a-and-all-source-independent-controls}{%
\section{4. Verify T03-A and all source-independent controls}\label{verify-t03-a-and-all-source-independent-controls}}

\begin{verbatim}
completed = subprocess.run(
    [sys.executable, str(RH_REPO / "tools" / "verify_t03.py")],
    cwd=RH_REPO,
    text=True,
    capture_output=True,
)
print(completed.stdout)
print(completed.stderr)
print("return code:", completed.returncode)
assert completed.returncode == 0
\end{verbatim}

Expected lines include:

\begin{verbatim}
T03-A PASS_T03_ACTION_DOMAIN_DENSE_CORE
T03-A certificate 05a4784308e996763928b4cead0c86462150154c506b30bf292f96dd8986db3c
T03-B OPEN: pinned private worktrees were not supplied
T03-C1 OPEN: actual finite action inventory was not extracted
\end{verbatim}

This reproduces T03-A and the fail-closed T03-B/T03-C1 control logic. It does not
execute the private MP packet.

\hypertarget{fetch-the-pinned-mp-branches}{%
\section{5. Fetch the pinned MP branches}\label{fetch-the-pinned-mp-branches}}

\begin{verbatim}
subprocess.run(
    [
        "git",
        "-C",
        str(MP_REPO),
        "fetch",
        "origin",
        "agent/native-recognition-operator-algebra",
        "agent/rh-rotation-metric-topology-reset",
    ],
    check=True,
)
\end{verbatim}

\hypertarget{create-clean-detached-source-worktrees}{%
\section{6. Create clean detached source worktrees}\label{create-clean-detached-source-worktrees}}

\begin{verbatim}
if not OPERATOR_WORKTREE.exists():
    subprocess.run(
        [
            "git",
            "-C",
            str(MP_REPO),
            "worktree",
            "add",
            "--detach",
            str(OPERATOR_WORKTREE),
            OPERATOR_COMMIT,
        ],
        check=True,
    )

if not STATE_WORKTREE.exists():
    subprocess.run(
        [
            "git",
            "-C",
            str(MP_REPO),
            "worktree",
            "add",
            "--detach",
            str(STATE_WORKTREE),
            STATE_COMMIT,
        ],
        check=True,
    )

print("operator head:", run_git(OPERATOR_WORKTREE, "rev-parse", "HEAD"))
print("state head:", run_git(STATE_WORKTREE, "rev-parse", "HEAD"))
print("operator dirty:", repr(run_git(OPERATOR_WORKTREE, "status", "--porcelain")))
print("state dirty:", repr(run_git(STATE_WORKTREE, "status", "--porcelain")))

assert run_git(OPERATOR_WORKTREE, "rev-parse", "HEAD") == OPERATOR_COMMIT
assert run_git(STATE_WORKTREE, "rev-parse", "HEAD") == STATE_COMMIT
assert run_git(OPERATOR_WORKTREE, "status", "--porcelain") == ""
assert run_git(STATE_WORKTREE, "status", "--porcelain") == ""
\end{verbatim}

The RH verifier rejects dirty worktrees, wrong commits, changed pinned blobs,
and altered label order.

\hypertarget{execute-the-pinned-carrier-and-t03-c1-inventory}{%
\section{7. Execute the pinned carrier and T03-C1 inventory}\label{execute-the-pinned-carrier-and-t03-c1-inventory}}

\begin{verbatim}
completed = subprocess.run(
    [
        sys.executable,
        str(RH_REPO / "tools" / "verify_t03.py"),
        "--operator-root",
        str(OPERATOR_WORKTREE),
        "--state-root",
        str(STATE_WORKTREE),
        "--t03-c1-output",
        str(T03_C1_JSON),
        "--t03-c1-npz-output",
        str(T03_C1_NPZ),
    ],
    cwd=RH_REPO,
    text=True,
    capture_output=True,
)
print(completed.stdout)
print(completed.stderr)
print("return code:", completed.returncode)
print("JSON certificate:", T03_C1_JSON)
print("NPZ arrays:", T03_C1_NPZ)
\end{verbatim}

Required status:

\begin{verbatim}
PASS_RH_FRAMEWORK_T03_C1_PINNED_ACTION_INVENTORY
\end{verbatim}

A passing run proves T03-B and T03-C1 together. A nonzero return code is a
verification failure, not a warning to be renamed until morale improves.

\hypertarget{inspect-theorem-states-and-false-gates}{%
\section{8. Inspect theorem states and false gates}\label{inspect-theorem-states-and-false-gates}}

\begin{verbatim}
import json

report = json.loads(T03_C1_JSON.read_text(encoding="utf-8"))
false_checks = [
    name for name, passed in report.get("checks", {}).items() if not passed
]
print("status:", report["status"])
print("classification:", report["theorem_classification"])
print("false checks:", false_checks)
print("certificate hash:", report.get("certificate_hash"))
\end{verbatim}

Required promotion state:

\begin{verbatim}
T03_B_finite_packet_reproduction = PROVED
T03_C1_actual_finite_packet_action_inventory = PROVED
T03_C_remaining_action_families = OPEN
riemann_hypothesis = OPEN
false checks = []
\end{verbatim}

\hypertarget{inspect-per-array-hashes-and-roles}{%
\section{9. Inspect per-array hashes and roles}\label{inspect-per-array-hashes-and-roles}}

\begin{verbatim}
inventory = report["action_inventory"]

print("carrier data:")
for name, metadata in inventory["carrier_data"].items():
    print(" ", name, metadata["shape"], metadata["sha256"])

print("matrix actions:")
for name, role in inventory["actions"].items():
    metadata = inventory["array_metadata"][name]
    print(
        " ",
        name,
        role["operator_class"],
        metadata["operator_norm"],
        metadata["sha256"],
    )
\end{verbatim}

The individual SHA-256 values identify the mathematical arrays. The NPZ file is
only their transport bundle.

\hypertarget{confirm-the-npz-names}{%
\section{10. Confirm the NPZ names}\label{confirm-the-npz-names}}

\begin{verbatim}
import numpy as np

with np.load(T03_C1_NPZ, allow_pickle=False) as loaded:
    print(sorted(loaded.files))
\end{verbatim}

Expected names include:

\begin{verbatim}
corrected_defect_line
source_defect_projector
accepted_complement_projector
source_metric
response_metric
loop_holonomy
memory_holonomy
compensated_holonomy
\end{verbatim}

\hypertarget{claim-boundary}{%
\section{Claim boundary}\label{claim-boundary}}

A passing run proves the pinned finite carrier reproduction and the actual
finite packet action inventory. T03-C2 through T03-C5 remain open: amplitude
families, arithmetic/Gamma actions, unbounded graph domains, and aperture
exhaustion. Target faithfulness, native Shunya, the critical-line adapter, and
RH remain open.

\begin{flushright}\small\texttt{RH-Framework/notebooks/T03\_ACTION\_DOMAIN\_JUPYTER.md}\end{flushright}
